\section{Worked computations ledger}

This chapter collects the exact computations used in the paper into a
single reference ledger.  Every entry is an exact integer identity.
The formal proof does not rely on this ledger as a decision procedure:
the corresponding facts are proved in Lean by explicit rewriting.

\subsection{Orbit and prefix ledger}

\begin{center}
\scriptsize
\begin{tabular}{@{}rrrrrrr@{}}
\toprule
$n$ & $\Orbit_7(n)$ & $t_n$ & $\wordWeight_n$ & $\wordA_n$ &
  $5^n\cdot7+\wordA_n$ & $2^{\wordWeight_n}\Orbit_7(n)$ \\
\midrule
0 & 7 & -- & 0 & 0 & 7 & 7 \\
1 & 9 & 2 & 2 & 1 & 36 & 36 \\
2 & 23 & 1 & 3 & 9 & 184 & 184 \\
3 & 29 & 2 & 5 & 53 & 928 & 928 \\
4 & 73 & 1 & 6 & 297 & 4672 & 4672 \\
5 & 183 & 1 & 7 & 1549 & 23424 & 23424 \\
6 & 229 & 2 & 9 & 7873 & 117248 & 117248 \\
7 & 573 & 1 & 10 & 39877 & 586752 & 586752 \\
8 & 1433 & 1 & 11 & 200409 & 2934784 & 2934784 \\
9 & 3583 & 1 & 12 & 1004093 & 14675968 & 14675968 \\
10 & 4479 & 2 & 14 & 5024561 & 73383936 & 73383936 \\
11 & 5599 & 2 & 16 & 25139189 & 366936064 & 366936064 \\
12 & 6999 & 2 & 18 & 125761481 & 1834745856 & 1834745856 \\
13 & 8749 & 2 & 20 & 629069549 & 9173991424 & 9173991424 \\
14 & 21873 & 1 & 21 & 3146396321 & 45871005696 & 45871005696 \\
15 & 54683 & 1 & 22 & 15734078757 & 229357125632 & 229357125632 \\
16 & 34177 & 3 & 25 & 78674588089 & 1146789822464 & 1146789822464 \\
\bottomrule
\end{tabular}
\end{center}

The last two columns are equal by the exact word-orbit identity
\begin{equation*}
  2^{\wordWeight_n}\,\Orbit_7(n)=5^n\cdot7+\wordA_n.
\end{equation*}
For example, the final row reads
\begin{equation*}
  2^{25}\cdot34177=5^{16}\cdot7+78674588089=1146789822464.
\end{equation*}

\subsection{Block equations}

\begin{align*}
  2^3\cdot23&=184=5^2\cdot7+9
  &&\text{block }[2,1]\text{ from }7,\\
  2^3\cdot29&=232=5^2\cdot9+7
  &&\text{block }[1,2]\text{ from }9,\\
  2^4\cdot229&=3664=5^3\cdot29+39
  &&\text{block }[1,1,2]\text{ from }29.
\end{align*}
The molecules are $\wordA([2,1])=9$, $\wordA([1,2])=7$, and
$\wordA([1,1,2])=39$.

\subsection{PMI sums}

For the exact prefix word $[2,1,2]$ from $7$,
\begin{equation*}
  \sum_{j=0}^{2}2^{-m_j}
    =1+\frac45+\frac8{25}
    =\frac{53}{25}
    =\frac{5\wordA(w)}{5^3}.
\end{equation*}
For the word $[3,1]$, the cleared sum is
\begin{equation*}
  \mathrm{aTotal5}(2,t)=5^2+2^3\cdot5^1=25+40=65=5\wordA([3,1]).
\end{equation*}
The first example has no bad prefix; the second has exactly one.

\subsection{Reduced segment equations}

\begin{align*}
  d=1:&& 5g+1&=2^{1+4a}\left(2^{e-1}g+\delta5^{j-1}\right),\\
  d=2:&& \left(2^{u_1+e}-25\right)g
    +2^{u_1+1}\delta5^{j-1}&=5+2^{u_1},\\
  d=3,\text{ survivor}:&& 16g+8\cdot5^{j-1}&=125g+39.
\end{align*}
The $d=3$ survivor therefore satisfies
\begin{equation*}
  g=\frac{8\cdot5^{j-1}-39}{109}.
\end{equation*}

\subsection{Modular tables}

Modulo $17$, the powers of $5$ satisfy
\begin{equation*}
  5^1\equiv5,\quad 5^2\equiv8,\quad 5^3\equiv6\pmod{17},
\end{equation*}
and the order is $16$.  Modulo $256$,
\begin{equation*}
  5^7=78125\equiv45\pmod{256},
\end{equation*}
and the order is $64$.  Modulo $109$, repeated squaring gives:
\begin{center}
\small
\begin{tabular}{@{}rrrrrrr@{}}
\toprule
$k$ & $5^k\bmod109$ & $k$ & $5^k\bmod109$ & $k$ & $5^k\bmod109$ \\
\midrule
1 & 5 & 16 & 89 & 256 & 26 \\
2 & 25 & 32 & 73 & 512 & 22 \\
4 & 80 & 64 & 97 & 923 & 73 \\
8 & 78 & 128 & 35 & & \\
\bottomrule
\end{tabular}
\end{center}
The entry $5^{923}\equiv73$ is the product
\begin{equation*}
  22\cdot26\cdot35\cdot89\cdot78\cdot25\cdot5
  \equiv73\pmod{109},
\end{equation*}
and $8\cdot73=584\equiv39\pmod{109}$, so the $d=3$ denominator is
consistent.

\subsection{Arithmetic counterexample}

The witness is
\begin{equation*}
  s=2^{83}-3\cdot5^{35}=940257419896922313665033.
\end{equation*}
The size comparison is
\begin{equation*}
  2^{81}=2417851639229258349412352
  <
  2910383045673370361328125=5^{35},
\end{equation*}
and the valuation identity is
\begin{equation*}
  5s+3\cdot5^{36}=5\cdot2^{83}=48357032784585166988247040.
\end{equation*}
Hence $\vtwo{5s+3\cdot5^{36}}=83$, while the unconstrained bound is
$82$.  The witness is odd, satisfies $s\equiv3\pmod5$, and fails only
the orbit-reachability fields.

\subsection{Verification note}

Every entry in this ledger is an exact integer identity.  The paper
uses these numbers only to illustrate the algebra; the formal proof
re-derives the corresponding facts by explicit rewriting in
\texttt{FinitePrefix.lean}, \texttt{WordWindow.lean}, and
\texttt{Pmi.lean}.  No entry here is a substitute for a Lean theorem.
